\chapter{2007年广东卷（理）}

\section{单选题}

\begin{problem}
已知函数 \( f(x) = \frac{1}{\sqrt{1 - x}} \) 的定义域为 \(M\)，\(g(x) = \ln (1 + x)\) 的定义域为 \(N\)，则 \( M \cap N = \)（\quad）

\choices
    {\( \{x \mid x > -1\} \)}
    {\( \{x \mid x < 1\} \)}
    {\( \{x \mid -1 < x < 1\} \)}
    {\(\varnothing\)}
\end{problem}

\begin{answer}
C
\end{answer}

\begin{solution}
函数 \(f(x)\) 有意义需要根式内为正且分母不为零，因此 \(1-x>0\)，即 \(x<1\)，所以 \(M=\{x\mid x<1\}\). 函数 \(g(x)=\ln(1+x)\) 有意义需要 \(1+x>0\)，所以 \(N=\{x\mid x>-1\}\). 因而
\[
M\cap N=\{x\mid -1<x<1\}
\]
故选 C.
\end{solution}

\begin{problem}
若复数 \((1 + b\i)(2 + \i)\) 是纯虚数（\(\i\) 是虚数单位，\(b\) 是实数），则 \(b =\)（\quad）

\choices
    {\(2\)}
    {\(\frac{1}{2}\)}
    {\(-\frac{1}{2}\)}
    {\(-2\)}
\end{problem}

\begin{answer}
A
\end{answer}

\begin{solution}
展开复数得
\[
(1+b\i)(2+\i)=2+\i+2b\i+b\i^2=(2-b)+(1+2b)\i
\]
纯虚数的实部必须为零，因此 \(2-b=0\)，解得 \(b=2\). 此时虚部为 \(1+2b=5\ne0\)，确为纯虚数，故选 A.
\end{solution}

\begin{problem}
若函数 \( f(x) = \sin^2 x - \frac{1}{2} \)（\(x \in \R\)），则 \( f(x) \) 是（\quad）

\choices
    {最小正周期为 \(\frac{\pi}{2}\) 的奇函数}
    {最小正周期为 \(\pi\) 的奇函数}
    {最小正周期为 \(2\pi\) 的偶函数}
    {最小正周期为 \(\pi\) 的偶函数}
\end{problem}

\begin{answer}
D
\end{answer}

\begin{solution}
利用降幂公式，
\[
f(x)=\sin^2x-\frac12=\frac{1-\cos2x}{2}-\frac12=-\frac12\cos2x
\]
由于 \(\cos2x\) 是偶函数，故 \(f(x)\) 是偶函数；\(\cos2x\) 的最小正周期为 \(\pi\)，且 \(f(x+\frac{\pi}{2})=-f(x)\) 不对所有 \(x\) 成立，所以 \(f(x)\) 的最小正周期不是 \(\frac{\pi}{2}\)，而是 \(\pi\). 故选 D.
\end{solution}

\begin{problem}
客车从甲地以 \(60 \, \mathrm{km/h}\) 的速度匀速行驶 \(1\) 小时到达乙地，在乙地停留了半小时，然后以 \(80 \, \mathrm{km/h}\) 的速度匀速行驶 \(1\) 小时到达丙地. 下列描述客车从甲地出发，经过乙地，最后到达丙地所经过的路程 \(s\) 与时间 \(t\) 之间关系的图象中，正确的是（\quad）

\choices
    {\choicebitmap[width=0.15\paperwidth]{img/2007/guangdong_science/q5_optA_clean.png}}
    {\choicebitmap[width=0.15\paperwidth]{img/2007/guangdong_science/q5_optB_clean.png}}
    {\choicebitmap[width=0.15\paperwidth]{img/2007/guangdong_science/q5_optC_clean.png}}
    {\choicebitmap[width=0.15\paperwidth]{img/2007/guangdong_science/q5_optD_clean.png}}
\end{problem}

\begin{answer}
B
\end{answer}

\begin{solution}
客车出发后前 \(1\) 小时行驶 \(60\) 千米，所以在 \(t=1\) 时 \(s=60\). 在乙地停留半小时，故 \(1\leqslant t\leqslant1.5\) 时路程保持为 \(60\). 随后以 \(80\,\mathrm{km/h}\) 行驶 \(1\) 小时，到达丙地时 \(t=2.5\)，总路程为 \(60+80=140\) 千米.

因此路程函数为
\[
s=\begin{cases}
60t,&0\leqslant t\leqslant1,\\
60,&1\leqslant t\leqslant1.5,\\
60+80(t-1.5),&1.5\leqslant t\leqslant2.5.
\end{cases}
\]
图象应先以斜率 \(60\) 上升，再水平停留半小时，最后以更大的斜率 \(80\) 上升至 \((2.5,140)\)，故选 B.
\end{solution}

\begin{problem}
已知数列 \(\{a_{n}\}\) 的前 \(n\) 项和 \(S_{n} = n^{2} - 9n\)，第 \(k\) 项满足 \(5 < a_{k} < 8\)，则 \(k =\)（\quad）

\choices
    {\(9\)}
    {\(8\)}
    {\(7\)}
    {\(6\)}
\end{problem}

\begin{answer}
B
\end{answer}

\begin{solution}
令 \(S_0=0\). 由前 \(n\) 项和得到
\[
a_k=S_k-S_{k-1}=k^2-9k-\bigl((k-1)^2-9(k-1)\bigr)=2k-10
\]
所以
\[
5<2k-10<8\quad\Longleftrightarrow\quad 15<2k<18\quad\Longleftrightarrow\quad \frac{15}{2}<k<9
\]
由于 \(k\) 是正整数，只能取 \(k=8\)，故选 B.
\end{solution}

\begin{problem}
图1是某县参加2007年高考的学生身高条形统计图，从左到右的各条形表示的学生人数依次记为 \(A_{1}\)，\(A_{2}\)，\(\dots\)，\(A_{10}\) （如 \(A_{2}\) 表示身高（单位：\(\mathrm{cm}\)）在 \([150, 155)\) 内的学生人数）. 图2是统计图1中身高在一定范围内学生人数的一个算法流程图. 现要统计身高在 \(160 \sim 180 \mathrm{~cm}\) （含 \(160 \mathrm{~cm}\)，不含 \(180 \mathrm{~cm}\)）的学生人数，那么在流程图中的判断框内应填写的条件是（\quad）

\bitmapfigure[float=inline,max width=0.9\linewidth]{img/2007/guangdong_science/q6_chart.png}
\texfigure[float=inline,max width=0.9\linewidth]{img_repaint/2007/guangdong_science/q6_flowchart.tex}

\choices
    {\(i < 6\)}
    {\(i < 7\)}
    {\(i < 8\)}
    {\(i < 9\)}
\end{problem}

\begin{answer}
C
\end{answer}

\begin{solution}
图1中 \(A_1\)，\(A_2\)，\(A_3\)，\(A_4\)，\(\ldots\) 分别对应区间 \([145,150)\)，\([150,155)\)，\([155,160)\)，\([160,165)\)，\(\ldots\). 因而身高在 \([160,180)\) 内的学生人数是
\[
A_4+A_5+A_6+A_7
\]
流程图先令 \(i=4\)，判断成立时累加 \(A_i\)，再令 \(i=i+1\). 为恰好累加 \(A_4\) 至 \(A_7\)，判断应在 \(i=4\)，\(5\)，\(6\)，\(7\) 时成立，在 \(i=8\) 时不成立，所以条件为 \(i<8\)，故选 C.
\end{solution}

\begin{problem}
如图是某汽车维修公司的维修点环形分布图. 公司在年初分配给 \(A\)，\(B\)，\(C\)，\(D\) 四个维修点某种配件各 \(50\) 件. 在使用前发现需将 \(A\)，\(B\)，\(C\)，\(D\) 四个维修点的这批配件分别调整为 \(40\)，\(45\)，\(54\)，\(61\) 件，但调整只能在相邻维修点之间进行，那么要完成上述调整，最少的调动件次（\(n\) 件配件从一个维修点调整到相邻维修点的调动件次为 \(n\)）为（\quad）

\FigureLayoutDeclare{img/2007/guangdong_science/q7_fig1.png}{3.5cm}{56bp 29bp 146bp 0bp}
\bitmapfigure[max width=0.30\linewidth]{img/2007/guangdong_science/q7_fig1.png}

\choices
    {\(15\)}
    {\(16\)}
    {\(17\)}
    {\(18\)}
\end{problem}

\begin{answer}
B
\end{answer}

\begin{solution}
把环形维修点按 \(A\to B\to C\to D\to A\) 编号. 相对于每点原有的 \(50\) 件，所需变化量分别为
\[
(-10,-5,4,11)
\]
设沿上述方向在四条边上的净调动量依次为 \(u\)，\(v\)，\(w\)，\(z\)，正值表示沿箭头方向调动. 各点的流入量减流出量给出
\[
\begin{cases}
z-u=-10\\
u-v=-5\\
v-w=4\\
w-z=11
\end{cases}
\]
令 \(u=t\)，解得 \(v=t+5\)，\(w=t+1\)，\(z=t-10\).
同一条边上若有相反方向的调动，可以相互抵消而不增加调动件次，因此总调动件次至少且可以取为
\[
\Phi(t)=|t|+|t+5|+|t+1|+|t-10|
\]
当 \(-1\leqslant t\leqslant0\) 时，
\[
\Phi(t)=-t+(t+5)+(t+1)+(10-t)=16
\]
四个绝对值的折点为 \(-5\)，\(-1\)，\(0\)，\(10\)，而 \(\Phi(t)\) 在相邻折点区间上的斜率依次为 \(-4\)，\(-2\)，\(0\)，\(2\)，\(4\)，所以 \(\Phi(t)\) 在 \([-1,0]\) 上取得最小值 \(16\)，故选 B.
\end{solution}

\begin{problem}
设 \(S\) 是至少含有两个元素的集合. 在 \(S\) 上定义了一个二元运算“*”（即对任意的 \(a\)，\(b \in S\)，对于有序元素对 \((a, b)\)，在 \(S\) 中有唯一确定的元素 \(a*b\) 与之对应）. 若对任意的 \(a\)，\(b \in S\)，有 \(a*(b*a)=b\)，则对任意的 \(a\)，\(b \in S\)，下列等式中不恒成立的是（\quad）

\choices
    {\((a*b)*a = a\)}
    {\([a * (b * a)] * (a * b) = a\)}
    {\( b*(b*b)=b \)}
    {\((a*b)*[b*(a*b)] = b\)}
\end{problem}

\begin{answer}
A
\end{answer}

\begin{solution}
已知运算满足
\[
a\ast(b\ast a)=b
\]
取其中的两个变量均为 \(b\)，立即得到
\[
b\ast(b\ast b)=b
\]
所以选项 C 恒成立.

由已知式直接把 \(a\ast(b\ast a)\) 替换为 \(b\)，再对 \(b\ast(a\ast b)\) 使用已知式（取第一变量为 \(b\)，第二变量为 \(a\)），得
\[
[a\ast(b\ast a)]\ast(a\ast b)=b\ast(a\ast b)=a
\]
所以选项 B 恒成立. 对已知式取第一变量为 \(a\ast b\)，第二变量为 \(b\)，得
\[
(a\ast b)\ast[b\ast(a\ast b)]=b
\]
所以选项 D 恒成立.

选项 A 不一定成立. 例如取 \(S=\{0,1\}\)，定义 \(u\ast v\) 为模 \(2\) 加法，则
\[
u\ast(v\ast u)=u+(v+u)\equiv v\pmod 2
\]
满足题设条件；但取 \(a=0\)，\(b=1\)，有 \((a\ast b)\ast a=1\ne0=a\). 故不恒成立的是 A.
\end{solution}

\section{填空题}

\begin{problem}
甲，乙两个袋中装有红，白两种颜色的小球，这些小球除颜色外完全相同. 其中甲袋装有\(4\)个红球，\(2\)个白球，乙袋装有\(1\)个红球，\(5\)个白球. 现分别从甲，乙两袋中各随机取出\(1\)个球，则取出的两球是红球的概率为\fillinblank{}. （答案用分数表示）
\end{problem}

\begin{answer}
\(\frac{1}{9}\)
\end{answer}

\begin{solution}
从甲袋取到红球的概率为 \(\frac{4}{6}=\frac23\)，从乙袋取到红球的概率为 \(\frac{1}{6}\). 两袋分别取球相互独立，因此取出的两球都是红球的概率为
\[
\frac23\cdot\frac16=\frac19
\]
\end{solution}

\begin{problem}
若向量 \(\bs{a}\)，\(\bs{b}\) 满足 \(|\bs{a}| = |\bs{b}| = 1\)，\(\bs{a}\)，\(\bs{b}\) 的夹角为 \(120^{\circ}\)，则 \(\bs{a} \cdot \bs{a} + \bs{a} \cdot \bs{b} = \)\fillinblank{}.
\end{problem}

\begin{answer}
\(\frac{1}{2}\)
\end{answer}

\begin{solution}
由数量积定义，\(\bs{a}\cdot\bs{a}=|\bs{a}|^2=1\)，且 \(\bs{a}\cdot\bs{b}=|\bs{a}|\,|\bs{b}|\cos120^{\circ}=-\frac12\).
所以
\[
\bs{a}\cdot\bs{a}+\bs{a}\cdot\bs{b}=1-\frac12=\frac12
\]
\end{solution}

\begin{problem}
在平面直角坐标系 \(xOy\) 中，有一定点 \(A(2,1)\). 若线段 \(OA\) 的垂直平分线过抛物线 \(y^{2} = 2px\)（\(p > 0\)）的焦点，则该抛物线的准线方程是\fillinblank{}.
\end{problem}

\begin{answer}
\(x=-\frac{5}{4}\)
\end{answer}

\begin{solution}
线段 \(OA\) 的斜率为 \(\frac12\)，中点为 \((1,\frac12)\)，所以其垂直平分线为
\[
y-\frac12=-2(x-1)
\]
即 \(y=-2x+\frac52\). 抛物线 \(y^2=2px\) 的焦点为 \((\frac p2,0)\). 将 \(y=0\) 代入垂直平分线，得到焦点横坐标 \(\frac54\)，故 \(\frac p2=\frac54\)，从而 \(p=\frac52\). 对于 \(y^2=2px\)，准线方程为 \(x=-\frac p2\)，因此准线为
\[
x=-\frac54
\]
\end{solution}

\begin{problem}
如果一个凸多面体是 \(n\) 棱锥，那么这个凸多面体的所有顶点所确定的直线共有\fillinblank{} 条. 这些直线中共有 \(f(n)\) 对异面直线，则 \(f(4) =\fillinblank{}\)； \(f(n) =\fillinblank{}\)（答案用数字或 \(n\) 的解析式表示）
\end{problem}

\begin{answer}
\(\frac{n(n+1)}{2}\)，\(12\)，\(\frac{n(n-1)(n-2)}{2}\)
\end{answer}

\begin{solution}
\(n\) 棱锥有 \(n\) 个底面顶点和一个顶点，共有 \(n+1\) 个顶点. 任取两个顶点确定一条直线，所以直线总数为
\[
\frac{(n+1)n}{2}
\]

任取两条由顶点确定的直线：若两条都是底面顶点连线，则它们共面；若两条都是侧棱，则它们相交于棱锥顶点. 因此异面直线恰由一条底面顶点连线和一条侧棱组成. 先取底面的一条直线，有 \(\frac{n(n-1)}2\) 种；再从不在这条底面直线上的其余 \(n-2\) 个底面顶点中选一个，与棱锥顶点连成侧棱. 此时两直线不相交，且不可能共面，因而恰为异面直线. 故
\[
f(n)=\frac{n(n-1)}2(n-2)=\frac{n(n-1)(n-2)}2
\]
代入 \(n=4\)，得 \(f(4)=12\).
\end{solution}

\section{选考题：填空题}

\begin{problem}
在平面直角坐标系 \(xOy\) 中，直线 \(l\) 的参数方程为 \(\begin{cases}
x = t + 3,\\
y = 3 - t,
\end{cases}\) （参数 \(t \in \R\)），圆 \(C\) 的参数方程为 \(\begin{cases}
x = 2\cos \theta,\\
y = 2\sin \theta + 2,
\end{cases}\) （参数 \(\theta \in [0, 2\pi]\)），则圆 \(C\) 的圆心坐标为\fillinblank{}，圆心到直线 \(l\) 的距离为\fillinblank{}.
\end{problem}

\begin{answer}
\((0,2)\)，\(2\sqrt{2}\)
\end{answer}

\begin{solution}
由圆的参数方程
\[
\begin{cases}
x=2\cos\theta\\
y=2\sin\theta+2
\end{cases}
\]
可知圆心为 \(C_0(0,2)\). 直线 \(l\) 的参数方程相加得
\[
x+y=(t+3)+(3-t)=6
\]
所以其方程为 \(x+y-6=0\). 圆心到直线的距离为
\[
\frac{|0+2-6|}{\sqrt{1^2+1^2}}=2\sqrt2
\]
\end{solution}

\begin{problem}
设函数 \( f(x) = |2x - 1| + x + 3 \)，则 \( f(-2) =\fillinblank{} \)；若 \( f(x) \leqslant 5 \)，则 \( x \) 的取值范围是\fillinblank{}.
\end{problem}

\begin{answer}
\(6\)，\([-1,1]\)
\end{answer}

\begin{solution}
\[
f(-2)=|2(-2)-1|+(-2)+3=5+1=6
\]
再分情况解不等式. 当 \(x<\frac12\) 时，
\[
f(x)=1-2x+x+3=4-x\leqslant5
\]
所以 \(x\geqslant-1\)，得到 \([-1,\frac12)\). 当 \(x\geqslant\frac12\) 时，
\[
f(x)=2x-1+x+3=3x+2\leqslant5
\]
所以 \(x\leqslant1\)，得到 \([\frac12,1]\). 合并两种情况，\(x\) 的取值范围为 \([-1,1]\).
\end{solution}

\begin{problem}
如图所示，圆 \(O\) 的直径 \(AB = 6\)，\(C\) 为圆周上一点，\(BC = 3\). 过点 \(C\) 作圆的切线 \(l\)，过 \(A\) 做 \(l\) 的垂线 \(AD\)，\(AD\) 分别与直线 \(l\)，圆交于点 \(D\)，\(E\)，则 \(\angle DAC =\)\fillinblank{}°，线段 \(AE\) 的长为\fillinblank{}.
\bitmapfigure[max width=0.34\linewidth]{img/2007/guangdong_science/q15_fig1.png}
\end{problem}

\begin{answer}
\(30\)，\(3\)
\end{answer}

\begin{solution}
取圆心 \(O\) 为原点，令 \(A(-3,0)\)，\(B(3,0)\). 设 \(C=(u,v)\) 且 \(v>0\)，由 \(C\) 在圆上及 \(BC=3\) 得
\[
\begin{cases}
u^2+v^2=9\\
(u-3)^2+v^2=9
\end{cases}
\]
两式相减得 \(u=\frac32\)，再代回得 \(v=\frac{3\sqrt3}{2}\). 因此 \(AC\) 的斜率为
\[
\frac{3\sqrt3/2}{3/2+3}=\frac1{\sqrt3}
\]
所以 \(AC\) 与 \(AB\) 的夹角为 \(30^{\circ}\).

半径 \(OC\) 的斜率为 \(\sqrt3\)，故切线 \(l\) 的斜率为 \(-\frac1{\sqrt3}\). 由于 \(AD\perp l\)，直线 \(AD\) 的斜率为 \(\sqrt3\)，它与 \(AB\) 的夹角为 \(60^{\circ}\). 因而
\[
\angle DAC=60^{\circ}-30^{\circ}=30^{\circ}
\]

直线 \(AD\) 的单位方向向量可取 \(\left(\frac12,\frac{\sqrt3}{2}\right)\). 设从 \(A\) 出发沿此方向到达圆的另一交点 \(E\) 的距离为 \(r>0\)，则
\[
E=A+r\left(\frac12,\frac{\sqrt3}{2}\right)
\]
由 \(E\) 在半径为 \(3\) 的圆上，
\(\left(-3+\frac r2\right)^2+\left(\frac{\sqrt3 r}{2}\right)^2=9\)，\(\Longrightarrow\)，\(r^2-3r=0\).
除去对应于 \(A\) 的 \(r=0\)，得 \(AE=r=3\).
\end{solution}

\section{解答题}

\begin{problem}
已知 \(\triangle ABC\) 顶点的直角坐标分别为 \(A(3,4)\)，\(B(0,0)\)，\(C(c,0)\).

\begin{enumerate}
\item 若 \(c = 5\)，求 \( \sin \angle A \) 的值；
\item 若 \(\angle A\) 是钝角，求 \(c\) 的取值范围.
\end{enumerate}
\end{problem}

\begin{solution}
\begin{enumerate}
\item \(\overrightarrow{AB}=(-3,-4)\)，\(\overrightarrow{AC}=(2,-4)\)，所以
\[
\sin\angle A=\frac{\left|\overrightarrow{AB}\times\overrightarrow{AC}\right|}{|\overrightarrow{AB}|\,|\overrightarrow{AC}|}=\frac{|(-3)(-4)-(-4)\cdot2|}{5\cdot2\sqrt5}=\frac{2\sqrt5}{5}
\]

\item \(\overrightarrow{AB}=(-3,-4)\)，\(\overrightarrow{AC}=(c-3,-4)\). 当且仅当
\[
\overrightarrow{AB}\cdot\overrightarrow{AC}=-3(c-3)+(-4)(-4)=25-3c<0
\]
时，\(\angle A\) 是钝角，因此 \(c>\dfrac{25}{3}\).
\end{enumerate}
\end{solution}


\begin{problem}
下表提供了某厂节能降耗技术改造后生产甲产品过程中记录的产量 \(x\) （吨）与相应的生产能耗 \(y\) （吨标准煤）的几组对照数据.

\begin{center}
\begin{tabular}{|c|c|c|c|c|}
\hline
\(x\) & \(3\) & \(4\) & \(5\) & \(6\) \\
\hline
\(y\) & \(2.5\) & \(3\) & \(4\) & \(4.5\) \\
\hline
\end{tabular}
\end{center}

\begin{enumerate}
\item 请画出上表的散点图；
\item 请根据上表提供的数据，用最小二乘法求出 \(y\) 关于 \(x\) 的线性回归方程 \( y = bx + \hat{a} \)；
\item 已知该厂技改前 \(100\) 吨甲产品的生产能耗为 \(90\) 吨标准煤. 试根据前一问求出的线性回归方程，预测生产 \(100\) 吨甲产品的生产能耗比技改前降低多少吨标准煤？（参考数值：\( 3 \times 2.5 + 4 \times 3 + 5 \times 4 + 6 \times 4.5 = 66.5 \)）
\end{enumerate}
\end{problem}

\begin{solution}
\begin{enumerate}
\item 散点图中的四个数据点为 \((3,2.5)\)，\((4,3)\)，\((5,4)\)，\((6,4.5)\). 在直角坐标系中描出这四个点，即得所求散点图.

\item 样本平均数为 \(\overline{x}=\frac{3+4+5+6}{4}=4.5\)，\(\overline{y}=\frac{2.5+3+4+4.5}{4}=3.5\). 由题给参考数值及数据表可得 \(\sum_{i=1}^{4}x_i y_i=66.5\)，\(\sum_{i=1}^{4}x_i^2=3^2+4^2+5^2+6^2=86\).
因此，最小二乘法所得回归直线的斜率为
\[
\widehat{b}=\frac{\sum x_i y_i-4\overline{x}\,\overline{y}}{\sum x_i^2-4\overline{x}^{\,2}}=\frac{66.5-4\cdot4.5\cdot3.5}{86-4\cdot4.5^2}=\frac{3.5}{5}=0.7
\]
截距为
\[
\widehat{a}=\overline{y}-\widehat{b}\,\overline{x}=3.5-0.7\cdot4.5=0.35
\]
所以线性回归方程为 \(\widehat{y}=0.7x+0.35\).

\item 当产量为 \(100\text{ 吨}\) 时，由回归方程预测生产能耗为
\[
\widehat{y}=0.7\times100+0.35=70.35\text{ 吨}
\]
因此比技改前降低
\[
90-70.35=19.65\text{ 吨}
\]
\end{enumerate}
\end{solution}

\begin{problem}
在平面直角坐标系 \(xOy\) 中，已知圆心在第二象限，半径为 \(2\sqrt{2}\) 的圆 \(C\) 与直线 \(y = x\) 相切于坐标原点 \(O\). 椭圆 \(\frac{x^2}{a^2} + \frac{y^2}{9} = 1\) 与圆 \(C\) 的一个交点到椭圆两焦点的距离之和为 \(10\).

\begin{enumerate}
\item 求圆 \(C\) 的方程；
\item 试探求 \(C\) 上是否存在异于原点的点 \(Q\)，使 \(Q\) 到椭圆右焦点 \(F\) 的距离等于线段 \(OF\) 的长. 若存在，请求出点 \(Q\) 的坐标. 若不存在，请说明理由.
\end{enumerate}
\end{problem}

\begin{solution}
\begin{enumerate}
\item 设圆心为 \(M(m,n)\). 因为圆在原点与直线 \(y=x\) 相切，所以半径 \(OM\) 垂直于直线 \(y=x\)，从而 \(m+n=0\). 又因为半径为 \(2\sqrt2\)，且原点在圆上，所以
\[
m^2+n^2=8
\]
由 \(m<0\)，\(n>0\)，得 \(M(-2,2)\)，故圆 \(C\) 的方程为
\[
(x+2)^2+(y-2)^2=8
\]

\item 椭圆上任一点到两焦点的距离之和为 \(2a\)，所以 \(2a=10\)，即 \(a=5\). 因而椭圆的焦半距为
\[
c=\sqrt{a^2-3^2}=\sqrt{25-9}=4
\]
椭圆右焦点为 \(F(4,0)\)，且 \(|OF|=4\).

若存在符合条件的点 \(Q(x,y)\)，则它同时满足
\[
\begin{cases}
(x+2)^2+(y-2)^2=8\\
(x-4)^2+y^2=16
\end{cases}
\]
两式分别化为 \(x^2+y^2+4x-4y=0\) 与 \(x^2+y^2-8x=0\)，相减得 \(y=3x\). 代入后得 \(10x^2-8x=0\)，所以 \(x=0\) 或 \(x=\frac45\).
当 \(x=0\) 时得到原点 \(O\)，舍去；当 \(x=\frac45\) 时，\(y=\frac{12}{5}\). 因此存在异于原点的点，且
\[
Q\left(\frac45,\frac{12}{5}\right)
\]
\end{enumerate}
\end{solution}

\begin{problem}
如图所示，等腰 \(\triangle ABC\) 的底边 \(AB = 6\sqrt{6}\)，高 \(CD = 3\)，点 \(E\) 是线段 \(BD\) 上异于点 \(B\)，\(D\) 的动点，点 \(F\) 在 \(BC\) 边上，且 \(EF \perp AB\). 现沿 \(EF\) 将 \(\triangle BEF\) 折起到 \(\triangle PEF\) 的位置，使 \(PE \perp AC\)，记 \(BE = x\)，\(V(x)\) 表示四棱锥 \(P - ACFE\) 的体积.

\begin{enumerate}
\item 求 \(V(x)\) 的表达式；
\item 当 \(x\) 为何值时，\(V(x)\) 取得最大值？
\item 当 \(V(x)\) 取得最大值时，求异面直线 \(AC\) 与 \(PF\) 所成的余弦值.
\end{enumerate}

\bitmapfigure[max width=0.34\linewidth]{img/2007/guangdong_science/q19_fig1.png}
\end{problem}

\begin{solution}
\begin{enumerate}
\item 建立空间直角坐标系，使 \(D\) 为原点，\(AB\) 在 \(xOy\) 平面内，取 \(A(-3\sqrt6,0,0)\)，\(B(3\sqrt6,0,0)\)，\(C(0,3,0)\).
因为 \(BE=x\)，且 \(E\) 在 \(BD\) 上，所以 \(E(3\sqrt6-x,0,0)\). 直线 \(BC\) 的斜率为 \(-\dfrac{1}{\sqrt6}\)，又 \(EF\perp AB\)，故 \(F\left(3\sqrt6-x,\frac{x}{\sqrt6},0\right)\)，\(EF=\frac{x}{\sqrt6}\)，且 \(0<x<3\sqrt6\).

折叠后，\(PE=BE=x\)，并且 \(PE\perp EF\). 由于 \(EF\) 的方向为 \(y\) 轴方向，设 \(P-E=(u,0,v)\). 又
\[
\overrightarrow{AC}=(3\sqrt6,3,0)
\]
且 \(PE\perp AC\)，所以 \(3\sqrt6u=0\)，即 \(u=0\). 由 \(PE=x\) 得 \(|v|=x\)，故四棱锥的高为 \(x\).

四边形 \(ACFE\) 与三角形 \(BEF\) 构成三角形 \(ABC\)，于是
\[
S_{ACFE}=S_{ABC}-S_{BEF}=\frac12\cdot6\sqrt6\cdot3-\frac12\cdot x\cdot\frac{x}{\sqrt6}=9\sqrt6-\frac{x^2}{2\sqrt6}
\]
因此
\[
V(x)=\frac13x\left(9\sqrt6-\frac{x^2}{2\sqrt6}\right)=\frac{x(108-x^2)}{6\sqrt6}
\]
且 \(0<x<3\sqrt6\).

\item 对上式求导，得
\[
V'(x)=\frac{36-x^2}{2\sqrt6}
\]
当 \(0<x<6\) 时，\(V'(x)>0\)；当 \(6<x<3\sqrt6\) 时，\(V'(x)<0\). 所以当 \(x=6\) 时，\(V(x)\) 取得最大值.

\item 当 \(x=6\) 时，
\(F\left(3\sqrt6-6,\sqrt6,0\right)\)，\(P\left(3\sqrt6-6,0,6\right)\).
取方向向量 \(\overrightarrow{AC}=(3\sqrt6,3,0)\)，\(\overrightarrow{PF}=(0,\sqrt6,-6)\)，则异面直线所成角的余弦为
\[
\cos\theta=\frac{|\overrightarrow{AC}\cdot\overrightarrow{PF}|}{|\overrightarrow{AC}|\,|\overrightarrow{PF}|}=\frac{3\sqrt6}{3\sqrt7\cdot\sqrt{42}}=\frac17
\]
\end{enumerate}
\end{solution}

\begin{problem}
已知 \(a\) 是实数，函数 \(f(x) = 2ax^2 + 2x - 3 - a\). 如果函数 \(y = f(x)\) 在区间 \([-1, 1]\) 上有零点，求 \(a\) 的取值范围.
\end{problem}

\begin{solution}
方程 \(f(x)=0\) 可改写为
\[
a(2x^2-1)=3-2x
\]
当 \(x=\pm\dfrac{1}{\sqrt2}\) 时，左边的系数 \(2x^2-1=0\)，而右边 \(3-2x\neq0\)，所以这两个点都不是方程的根. 对其余的 \(x\in[-1,1]\)，有
\[
a=g(x)=\frac{3-2x}{2x^2-1}
\]
又
\[
g'(x)=\frac{4x^2-12x+2}{(2x^2-1)^2}
\]
在区间 \(\left[-1,-\dfrac{1}{\sqrt2}\right)\) 上，\(g\) 递增，且 \(g(-1)=5\)，\(\displaystyle\lim_{x\to-1/\sqrt2^-}g(x)=+\infty\)，故此时 \(g(x)\) 的取值范围为 \([5,+\infty)\).

在区间 \(\left(-\dfrac{1}{\sqrt2},\dfrac{1}{\sqrt2}\right)\) 上，\(g\) 在 \(x_0=\dfrac{3-\sqrt7}{2}\) 处取得最大值，并且两端的极限均为 \(-\infty\). 由
\[
g(x_0)=\frac{\sqrt7}{7-3\sqrt7}=-\frac{3+\sqrt7}{2}
\]
可知此时的取值范围为 \(\left(-\infty,-\dfrac{3+\sqrt7}{2}\right]\).

在区间 \(\left(\dfrac{1}{\sqrt2},1\right]\) 上，\(g\) 递减，且 \(g(1)=1\)，\(\displaystyle\lim_{x\to1/\sqrt2^+}g(x)=+\infty\)，故此时的取值范围为 \([1,+\infty)\). 合并上述范围，得
\[
a\in\left(-\infty,-\frac{3+\sqrt7}{2}\right]\cup[1,+\infty)
\]
\end{solution}

\begin{problem}
已知函数 \( f(x) = x^2 + x - 1 \)，\(\alpha\)，\(\beta\) 是方程 \( f(x) = 0 \) 的两个根（\( \alpha > \beta \)），\( f'(x) \) 是 \( f(x) \) 的导数. 设 \(a_1 = 1\)，\(a_{n+1} = a_n - \frac{f(a_n)}{f'(a_n)}\)（\(n = 1\)，\(2\)，\(\cdots\)）.

\begin{enumerate}
\item 求 \(\alpha\)，\(\beta\) 的值；
\item 证明：对任意的正整数 \(n\)，都有 \(a_{n} > \alpha\)；
\item 记 \( b_{n} = \ln \frac{a_{n} - \beta}{a_{n} - \alpha} \)（\(n = 1\)，\(2\)，\(\cdots\)）. 求数列 \(\{b_{n}\}\) 的前 \( n \) 项和 \( S_{n} \).
\end{enumerate}
\end{problem}

\begin{solution}
\begin{enumerate}
\item 由求根公式，得
\(\alpha=\frac{-1+\sqrt5}{2}\)，\(\beta=\frac{-1-\sqrt5}{2}\).

\item 因为 \(f'(x)=2x+1\)，递推式可化为
\[
a_{n+1}=a_n-\frac{a_n^2+a_n-1}{2a_n+1}=\frac{a_n^2+1}{2a_n+1}
\]
初始时 \(a_1=1>\alpha\). 假设 \(a_n>\alpha\)，则 \(2a_n+1>0\)，并利用 \(\alpha^2+\alpha-1=0\)，有
\[
a_{n+1}-\alpha=\frac{a_n^2+1-\alpha(2a_n+1)}{2a_n+1}=\frac{(a_n-\alpha)^2}{2a_n+1}>0
\]
所以 \(a_{n+1}>\alpha\). 由数学归纳法，对任意正整数 \(n\)，都有 \(a_n>\alpha\).

\item 由递推式，
\[
a_{n+1}-\beta
=\frac{a_n^2+1-\beta(2a_n+1)}{2a_n+1}
\]
又因 \(\beta^2+\beta-1=0\)，即 \(1-\beta=\beta^2\)，所以
\[
a_{n+1}-\beta
=\frac{a_n^2-2\beta a_n+1-\beta}{2a_n+1}
=\frac{a_n^2-2\beta a_n+\beta^2}{2a_n+1}
=\frac{(a_n-\beta)^2}{2a_n+1}
\]
由于 \(a_n>\alpha>\beta\)，两边均为正数，于是
\[
\frac{a_{n+1}-\beta}{a_{n+1}-\alpha}=\left(\frac{a_n-\beta}{a_n-\alpha}\right)^2
\]
取对数，得到 \(b_{n+1}=2b_n\)，所以 \(\{b_n\}\) 是首项为
\[
b_1=\ln\frac{1-\beta}{1-\alpha}=\ln\frac{3+\sqrt5}{3-\sqrt5}
\]
且公比为 \(2\) 的等比数列. 因此
\[
S_n=b_1\left(1+2+\cdots+2^{n-1}\right)=(2^n-1)\ln\frac{3+\sqrt5}{3-\sqrt5}
\]
\end{enumerate}
\end{solution}
